\documentclass[10pt,a4paper]{amsart}
\usepackage[margin=19mm]{geometry}
\usepackage{amsmath,amssymb,mathtools}
\usepackage[scr=rsfs,cal=euler]{mathalfa}
\usepackage{xfrac}
\usepackage[unicode,hidelinks,bookmarks=false]{hyperref}
\newcommand{\ie}{i.e.\ }
\newcommand{\cf}{cf.\ }
\newcommand{\resp}{respectively\ }
\newcommand{\Subsection}{\subsection}
\providecommand{\nobreakdash}{\nobreak\hbox{-}}
\setlength{\emergencystretch}{2em}
\multlinegap0pt
\usepackage{float}
\usepackage{amssymb}
\newtheorem{prop}{Proposition}
\newtheorem{lemma}{Lemma}
\newcommand{\bbR}{\mathbb{R}}
\newcommand{\frk}{\mathfrak{k}}
\newcommand{\frt}{\mathfrak{t}}
\title{Distribution of spin norm along pencils: the $Sp(p, q)$ case}
\author{Chao-Ping Dong, Zhan Ying}
\date{Source: arXiv:2605.06015v1 (CC BY 4.0)}
\begin{document}
\maketitle

\noindent\emph{Source and licence.} Distribution of spin norm along pencils: the $Sp(p, q)$ case. Version arXiv:2605.06015v1 (CC BY 4.0), distributed by arXiv under the Creative Commons Attribution 4.0 International licence, http://creativecommons.org/licenses/by/4.0/, as recorded in the arXiv licence metadata for this exact version and shown on the abstract page https://arxiv.org/abs/2605.06015v1. Full licence notice: \texttt{licenses/arxiv-2605.06015-CC-BY-4.0.txt}.\par\smallskip

\noindent\textit{Editorial scope.} Counted unit: the article's lemma lemma-down (source0.tex lines 590--601). The article's complete original proof of it is delivered with it (source0.tex lines 602--682, one \texttt{proof} environment of 81 lines). No mathematics is written, repaired or re-derived: the statement and the proof are the article's own lines, in the article's own notation, and no retained line is substituted or re-flowed.  Retained uncounted as the source-local context the proof needs: lines 295--297; lines 303--305; lines 338--340; lines 416--417; lines 510--516.  Nothing was omitted from any retained span, so the package carries no omission record.  No retained line contains a citation, so the wrapper carries no bibliography and no bibliography entry.  No retained line contains a figure environment, an image-inclusion command or drawing code, so no image is delivered and the package declares no asset dependency.  Editorial formatting only, in addition: an AMS \texttt{amsart} preamble that restates the article's own declarations and macros used by the retained lines, namely the environments \texttt{prop}, \texttt{lemma} on separate counters, together with the standard packages those lines need.  The retained environments print in this document as Proposition 1, Lemma 1 in order of appearance, whereas the article prints the same items with the numbering of its own full document.  The complete native source of the article as distributed in the arXiv e-print for this exact version is bundled unchanged as \texttt{sources/199992/source0.tex}.
\begin{equation}\label{mu}
\mu=(a_1,\dots,a_p\mid b_1,\dots,b_q)
 \end{equation}

\begin{equation}\label{eq4}
\sum_{i=1}^f a_i+\sum_{j=1}^g b_j>2pq-2(p-f)(q-g).
\end{equation}

\begin{equation}\label{sum-mufg}
S(\mu_{f, g})> 2pq.
\end{equation}

\begin{equation}\label{eq5}
\rho_n^{(\ell)}-\rho_n^{(L_{\ell, j})}=(\underbrace{0,\dots,0,1}_{j},0,\dots,0\mid 0,\dots,0,\underbrace{-1,0,\dots,0}_{k_{j}}).\end{equation}

\begin{prop}\label{prop-2p+1}
	For any u-large $\mu$ in \eqref{mu} such that $\mu-\beta$ is $\Delta^+(\frk,\frt)$-dominant, we have
	\begin{itemize}
		\item[(a)] If $b_1\ge 2p+1$, then $\Vert\mu\Vert_{\mathrm{spin}}>\Vert\mu-\beta\Vert_{\mathrm{spin}}$.
		\item[(b)] If $a_1\ge 2q+1$, then $\Vert\mu\Vert_{\mathrm{spin}}>\Vert\mu-\beta\Vert_{\mathrm{spin}}$.
	\end{itemize}
\end{prop}

\begin{lemma}\label{lemma-down}
Take $\mu\in \Lambda_{f, g}$ for certain $1\leq f\leq p$ and $1\leq g\leq q$.
Let $\rho_n^{(\ell)}=(k_1,\dots,k_p\mid r_1,\dots,r_q)\in\Omega_{p,q}$. Assume that $k_1>a_1$.
	\begin{itemize}
		\item[(a)] If $\mu$ is a $\partial R$-weight, then exists  $\rho_n^{(\ell\downarrow)}=(K_{\ell\downarrow}\mid R_{\ell\downarrow})\in\Omega_{p, q}$ such that $\{\mu_\ell\}\gg\{\mu_{\ell\downarrow}\}$ and that
		\[K_{\ell\downarrow}=(a_1, \underbrace{a_1-1, \dots, a_1-1}_{h-1}, k_{h+1}, \dots, k_p),\]
		where $h$ is the largest index satisfying $k_h\ge a_1$.
		\item[(b)] If $\mu$ is a $\partial L$-weight, then exists $\rho_n^{(\ell\downarrow)}=(K_{\ell\downarrow}\mid R_{\ell\downarrow})\in\Omega_{p, q}$ such that $\{\mu_\ell\}\gg\{\mu_{\ell\downarrow}\}$ and that
		\[R_{\ell\downarrow}=(b_1, \underbrace{b_1-1, \dots, b_1-1}_{h-1}, r_{h+1}, \dots, r_q),\]
		where $h$ is the largest index satisfying $r_h\ge b_1$.
	\end{itemize}
\end{lemma}

\begin{proof}
	We  only prove (a).  Note that $a_1=1$ can not happen. Indeed, we would then have
	\[2pg+2f(q-g)=2pq-2(q-f)(p-g)<\sum_{i=1}^f a_i+\sum_{j=1}^g b_j\le 2+(2p-1)g.\]
  It follows that $g+2f(q-g)<2$. This inequality holds if and only if $g\le 1$ and $f=0$, contradicting the assumption. Therefore, $a_1\ge 2$.


Firstly, take $\mu\in\Lambda_{p, q}$. Then
	\[\sum_{i=1}^pa_i+\sum_{j=1}^qb_j>2pq.\]
	Let $a$ be the minimum of the first components among all $\partial R$-weights in $\Lambda_{p,q}$, then $a\ge 2$. Denote
	\begin{align*}
		\nu&=(a,a-1,\dots,a-1\mid 2p,2p-1,\dots,2p-1), \\
		\nu^*&=(a-1,a-2,\dots,a-2\mid 2p,2p-1,\dots,2p-1).
	\end{align*}
	It is clear that $\nu$ is in $\Lambda_{p,q}$ and $\nu^*$ is not in $\Lambda_{p,q}$. Moreover, $S(\nu)-S(\nu^*)=p$.
Take any $\partial R$-weight $\kappa$ in $\Lambda_{p,q}$ whose first component is $a$. Then
$$
S(\nu^*)\le 2pq< S(\kappa)\leq S(\nu)
$$
If the last component of $\kappa$ is less than or equal to $p-1$, then
$$
S(\kappa)\leq S(\nu)-p=S(\nu)-\big(S(\nu)-S(\nu^*)\big)=S(\nu^*).
$$
This is a contradiction. Thus the last component of $\kappa$ is at least $p$. Hence every $\partial R$-weight in $\Lambda_{p,q}$ with first component $a$ has all of its last $q$ components at least $p$. Furthermore, denote
	\[\widetilde{\nu}=(a+1,a,\dots,a\mid 2p,2p-1,\dots,2p-1).\]
	Then $S(\widetilde{\nu})-S(\nu)=p$.
Take any $\partial R$-weight $\widetilde{\kappa}$ in $\Lambda_{p,q}$ with first component $a+1$. Then
$$
S(\nu^*)<S(\widetilde{\kappa})\leq S(\widetilde{\nu}).
 $$
If the $(n-1)$-th component of $\kappa$ is less than or equal to $p-1$, then
$$
S(\widetilde{\kappa})\leq S(\widetilde{\nu})-2p=S(\nu)-p\leq S(\nu^*).
$$
This is a contradiction. Hence the $(n-1)$-th component of any $\partial R$-weight in $\Lambda_{p,q}$ with first component $a+1$ is at least $p$. This procedure continues until $a$ reaches $a_1$,  which is the first component of $\mu$. When the procedure terminates, we see that $b_{q-(a_1-a)}$ is not less than $p$. Thus $b_{q-a_1+1}\ge b_{q-(a_1-a)}\ge p$.
	
	Since \eqref{eq5} shows that
	\[\rho_n^{(\ell)}-\rho_n^{(L_{\ell, h})}=(\underbrace{0,\dots,0,1}_{h},0,\dots,0\mid 0,\dots,0,\underbrace{-1,0,\dots,0}_{k_{h}}).\]
	Note that the $h$-th component $a_h-k_h$ of $\mu_\ell$ is a negative integer. Moreover, $k_h>0$ implies that the $(n-k_h+1)$-th component of $\rho_n^{(\ell)}$ is less than $p$. Hence, by the above argument and $n-k_h+1\le n-a_1+1$, the $(n-k_h+1)$-th component of $\mu$ is at least $p$. Thus the $(n-k_h+1)$-th component of $\mu_\ell$ is a positive integer. Now, the same argument as in Proposition \ref{prop-2p+1} shows that after a finite number of operations, we obtain the desired result.
	
	Now we turn to the general case. That is, $\mu\in \Lambda_{f, g}$. We \emph{claim} that $q-a_1+1\le g$. It is obviously true if $g=q$. Now assume that $g<q$.
	Indeed, if $g \le q/2$, then \eqref{eq4} shows that
	\[\sum_{i=1}^f a_i+\sum_{j=1}^g b_j>2pq-2(p-f)(q-g) = 2pg+2f(q-g)\ge 2pg+fq.\]
On the other hand,
$$
\sum_{i=1}^f a_i+\sum_{j=1}^g b_j\leq fq+2pg
$$
since each $a_i$ is at most $q$ and each $b_j$ is at most $2p$. This is a contradiction. Therefore, $g>q/2$.
	
	Assume that $3q/4\ge g>q/2$. Then $a_1\le q/2$ can not happen. Indeed, we would then have
	\[\dfrac{1}{2}fq+2pg \ge \sum_{i=1}^f a_i+\sum_{j=1}^g b_j >2pg+2f(q-g).\]
	This implies that $g>\dfrac{3}{4}q$, which is a contradiction. Hence $a_1>q/2$. Then
	\[q-a_1+1\leq\left[\dfrac{1}{2}q \right]+1\le g.\]
	Here for any $x\in\bbR$, $[x]$ denotes the smallest integer which is less than or equal to $x$. Thus the claim holds in this case.
	
	More generally, assume that $(1-1/2^{d+1})q\ge g>(1-1/2^d)q$. Then $a_1\le q/2^d$ can not happen. Indeed, we would then have
	\[\dfrac{1}{2^{d}}fq+2pg\ge \sum_{i=1}^f a_i+\sum_{j=1}^g b_j>2pg+2f(q-g).\]
	This implies that
	$g>(1-1/2^{d+1})q$, which is a contradiction.
	Therefore, $a_1>q/2^{d}$, then
	\[q-a_1+1\le \left[(1-\dfrac{1}{2^d})q \right]+1\le g.\]
	Thus the claim holds in this case.

Since $q$ is a fixed constant, there exists a positive integer $D$ such that $2^D\ge q$. Then
$$
(1-1/2^D)q\ge q-1.
$$
Therefore, the above procedure terminates when $d=D$ and the claim follows. It follows from the claim that $b_{q-a_1+1}\ge b_g$.
	
	Now let $a\ge 2$ be the minimum of the first components of all $\partial R$-weights in $\Lambda_{f,g}$. Put
	\[\nu=(\underbrace{a,a-1,\dots,a-1}_{f},a-2,a-2,\dots,a-2 \mid 2p,2p-1,\dots,2p-1),\]
	then $\nu$ is a $\partial R$-weight in $\Lambda_{f,g}$. Thus $S(\nu_{f, g})>2pq$ by \eqref{sum-mufg}. Construct
	\[\nu^*=(\underbrace{a-1,a-2,\dots,a-2}_{f},a-2,a-2,\dots,a-2 \mid 2p,2p-1,\dots,2p-1),\]
	then $\nu^*$ not in $\Lambda_{f,g}$. Thus $S(\nu_{f, g}^*)\leq 2pq$ by \eqref{sum-mufg}. Note that $S(\nu_{f,g})-S(\nu_{f,g}^*) = f\le p$. Take any $\partial R$-weight $\kappa$ in $\Lambda_{f,g}$ whose first component is $a$. Then
	\[S(\nu^*_{f,g}) < S(\kappa_{f,g}) \le S(\nu_{f,g}).\]
If the $(p+g)$-th component of $\kappa$ is less than or equal to $p-1$, then
\[S(\kappa_{f, g})\leq S(\nu_{f,g})-p\leq S(\nu_{f,g})-f=S(\nu_{f,g})-\big(S(\nu_{f,g})-S(\nu_{f,g}^*)\big)=S(\nu_{f,g}^*).\]
This is a contradiction. Hence the $(p+g)$-th component of any $\partial R$-weight in $\Lambda_{f,g}$ with first component $a$ is at least $p$. In other words, $b_1\ge \dots \ge b_g\ge p$.  Notice that
	\[a_1-a\le q-a< g,\]
where the last inequality follows from the claim.
Proceeding as in the case of $\mu\in \Lambda_{p, q}$, we obtian that $b_{g-(a_1-a)}\ge p$. Thus $b_{g-a_1+1}\ge b_{g-(a_1-a)}\ge p$. Therefore the discussion in the $\Lambda_{p, q}$ case applies here.
\end{proof}

\end{document}
